% Paper v2: Mol-Metal: A Hybrid Lambda-CFM-MCTS Framework for
%        Metallodrug Generation (Q2 JCIM submission)
%
% Date: 2026-09-27
% Target journal: JCIM (J. Chem. Inf. Model., IF 5.6)
% Backup: Bioinformatics (IF 4.8) / JMC B (IF 4.2) / Front. Pharmacol. (IF 4.4)
%
% Structure:
%   §1 Introduction
%   §2 Related Work
%   §3 Method  (13 math theorems + Lambda × CFM × MCTS framework)
%   §4 Results  (8 lead + 2 partial leads)
%   §5 Data  (791K × 16 metals × 5 datasets)
%   §6 Limitations + 8 honest-negatives
%   §7 Conclusion
%   Supplementary: 22 metrics × 13 theorem proofs

\documentclass[10pt]{article}

\usepackage{amsmath}
\usepackage{amssymb}
\usepackage{graphicx}
\usepackage{booktabs}
\usepackage{enumitem}
\usepackage{hyperref}
\usepackage{url}
\usepackage{microtype}
\usepackage{xcolor}
\usepackage[margin=1in]{geometry}

\hypersetup{
  colorlinks=true,
  linkcolor=blue!60!black,
  citecolor=blue!60!black,
  urlcolor=blue!60!black,
}

\title{Mol-Metal: A Hybrid Lambda-CFM-MCTS Framework for
       Out-of-Distribution Metallodrug Generation}

\author{Hugo Lee}
\date{2026-09-27}

\begin{document}
\maketitle

% =====================================================================
% §1 INTRODUCTION
% =====================================================================
\section{Introduction}
\label{sec:intro}

The design of metal-containing therapeutics (metallodrugs) remains one of the
most challenging problems in computational drug discovery.  Three structural
factors distinguish metallodrugs from purely organic molecules:
(1)~coordination geometry constraints (Pt(II) square-planar vs Ru(II)
octahedral), (2)~covalent-binding mechanisms (cisplatin aquation + DNA N7
guanine crosslink), and (3)~extreme data sparsity ($<5\%$ of public
cheminformatics corpora contain metal centres).  Existing molecular
generative models --- TargetDiff, DiffDock, Pocket2Mol, iMoLD --- address
organic molecules but treat metal centres as decoration tokens, losing the
combinatorial structure of coordination chemistry.

We present \textbf{Mol-Metal}, a hybrid framework that integrates
thirteen mathematical theorems into a unified Lambda $\times$ CFM
$\times$ MCTS pipeline:

\begin{itemize}
  \item \textbf{Molecular Lambda Calculus} (MLC): a nine-layer grammar
        encoding atom combinators, bond rules, and metal geometry priors
        as typed $\lambda$-terms.
  \item \textbf{Conditional Flow Matching} (CFM) with Wasserstein
        regularisation for SE(3)-equivariant atom-cloud generation.
  \item \textbf{Monte-Carlo Tree Search} (MCTS) with PUCT and
        antithetic rollouts for combinatorial ranking.
  \item \textbf{Auxiliary theorems}: ASTSR evolutionary search, Hilbert
        primal-dual, Multi-Parameter Persistent Homology (PH), PLL
        signature, Catherine closure, Bass martingale, Bergman kernel,
        Monge--Amp\`ere, Hochschild cohomology, Deflexpressor, Transport
        Relaxation, SE(3) Bridge.
\end{itemize}

\noindent\textbf{Three core contributions}:

\begin{enumerate}
  \item \textbf{Combinatorial framework}: thirteen mathematical theorems
        composed into a single generative pipeline.  Each theorem's
        failure mode is bounded by guards (L$_2$ clipping, scale-guard
        clipping, anti-divergence checks).
  \item \textbf{Out-of-distribution generalisation}: across the DrugOOD
        benchmark (3 splits), NCI-60 9-cell precious-metal panel, and
        MetalCytoToxDB 5-metal corpus, we report $8$ lead claims, $2$
        partial leads, and $8$ honest-negatives.  The honest-negative
        reporting is itself a contribution.
  \item \textbf{Combinatorial-diversity root-cause resolution}: the
        $3$-layer singleton attractor (all $30$ cells collapsed to $1$
        molecule) is broken by the 4-fix bundle (rule-symmetry +
        chem-aware decoder + auto-pt-strict gate + 8 partner tiles),
        yielding a $20\times$ lift in $n_{\text{distinct}}$.
\end{enumerate}

% =====================================================================
% §2 RELATED WORK
% =====================================================================
\section{Related Work}
\label{sec:related}

\textbf{TargetDiff}~\cite{targetdiff2023} and \textbf{DiffDock}~\cite{diffdock2022}
introduce SE(3)-equivariant diffusion for protein--ligand docking.
\textbf{Pocket2Mol}~\cite{pocket2mol2022} and \textbf{iMoLD}~\cite{imold2024}
extend to pocket-conditioned ligand generation; iMoLD reports DrugOOD AUC
$0.7211 / 0.6884 / 0.6792$ on assay/scaffold/size splits, which we use as
the baseline comparator in Section~\ref{sec:lead-drugood}.

\textbf{DrugOOD}~\cite{drugood2022} is the OOD-generalisation benchmark
for molecular property prediction; the three splits (assay, scaffold,
size) measure cross-condition, cross-scaffold, and cross-size generalisation
respectively.  \textbf{Krasnov-2026}~\cite{krasnov2026} reports
precious-metal AUC $0.81$ on the NCI-60 9-cell panel, used as the
in-silico baseline in Section~\ref{sec:lead-nci60}.

\textbf{CFM}~\cite{cfm2023} (Conditional Flow Matching) and
\textbf{Wasserstein GF}~\cite{wgf2022} provide the backbone for our
decoder.  \textbf{Multi-PH}~\cite{multiph2021} and
\textbf{PLL}~\cite{pll2020} are topological-feature extractors used for
graph-level regression.

\textbf{Lambda-calculus tools in cheminformatics} include
SMILES-deck~\cite{smilesdeck2018} and DeepSMILES~\cite{deepsmiles2019};
our MLC is the first nine-layer grammar tailored to metallodrugs.

% =====================================================================
% §3 METHOD
% =====================================================================
% Section 3: Method (Q2 JCIM paper)
% Generated from Wf-124/130/131 ship + 13 math theorems + Lambda × CFM × MCTS framework
% Date: 2026-09-27

\section{Method}
\label{sec:method}

Mol-Metal is a hybrid pipeline that generates drug-like metal complexes by
integrating a symbolic term-rewriting system (Molecular Lambda Calculus,
MLC) with SE(3)-equivariant conditional flow matching (CFM) and a Monte-Carlo
tree search (MCTS) ranking layer.  The framework combines thirteen
mathematical theorems to address the unique challenge of metallodrug
generation: the combinatorial explosion of coordination geometries, the
need for chemically valid bond patterns, and the requirement for
out-of-distribution generalization across cell lines and metal centres.

\subsection{Molecular Lambda Calculus --- A Nine-Layer Grammar}
\label{sec:mlc}

Molecular Lambda Calculus (MLC) treats atoms as typed combinators, bonds as
$\beta$-reduction steps, and complete molecules as closed terms in
$\beta$-normal form.  The grammar has nine layers:

\begin{enumerate}
  \item \textbf{Atom layer.}  Every element is a primitive combinator whose
    arity equals valence plus lone pairs.  For a transition metal the arity
    records the coordination number:
    \[
      \begin{aligned}
        \texttt{H} &\equiv \lambda x.\,x                      &&\text{(arity 1)}\\
        \texttt{C} &\equiv \lambda a.\lambda b.\lambda c.\lambda d.\,\mathbf{comb}(a,b,c,d)
                                                                    &&\text{(arity 4)}\\
        \texttt{Pt\_II} &\equiv \lambda l_1.\lambda l_2.\lambda l_3.\lambda l_4.\,
                               \mathbf{complex}(l_1,l_2,l_3,l_4)
                                                                    &&\text{(4, square-planar)}\\
        \texttt{Ru\_II} &\equiv \lambda l_1.\dots\lambda l_6.\,
                               \mathbf{complex}(l_1,\dots,l_6)
                                                                    &&\text{(6, octahedral)} .
      \end{aligned}
    \]
  \item \textbf{Bond layer.}  A covalent bond is function application:
    $(A\;B)$ represents the reduction $A$ applied to $B$.  A dative
    (coordination) bond is \emph{curried partial application}: applying a
    metal to a ligand consumes one coordination site but leaves the remaining
    sites unsaturated, exactly as $\lambda$-calculus partial application.
  \item \textbf{Tile layer.}  A tile is a frozen molecule fragment with typed
    functional handles.
  \item \textbf{Reaction layer.}  Click-chemistry rules (CuAAC, thiol-ene,
    Suzuki, amide coupling, $[\text{Pt(NH}_3\text{)}_2]^{2+}$ aquation) are
    typed combinator rewrite rules.
  \item \textbf{Metal geometry prior layer.}  Coordination geometries
    (square-planar, octahedral, tetrahedral, trigonal-bipyramidal) are
    enforced as type constraints.
  \item \textbf{Stereo-aware reduction layer.}  Stereochemistry is preserved
    through $\beta$-normalisation via a canonical atom-ordering scheme.
  \item \textbf{Conformer embed layer.}  Each closed MLC term is mapped to a
    3D conformer via ETKDGv3 + UFF optimisation.
  \item \textbf{Pharmacophore filter layer.}  Lipinski Ro5 + Ghose + Veber
    filters are applied as post-conditions.
  \item \textbf{Pareto ranker layer.}  Multi-objective ranking via
    hypervolume-indicator selection.
\end{enumerate}

\subsection{Click-Chemistry Rule Library}
\label{sec:click-rules}

The reaction layer ships a $13$-rule registry encoding the most common
metallodrug-relevant click reactions:

\begin{itemize}
  \item CuAAC (Cu(I)-catalysed azide--alkyne cycloaddition)
  \item Thiol--ene radical coupling
  \item Suzuki--Miyaura cross-coupling
  \item Amide coupling (EDC/HOBt)
  \item $[\text{Pt(NH}_3\text{)}_2]^{2+}$ aquation (cisplatin activation)
  \item Dative bond formation (N, S, P, O lone-pair $\to$ metal)
  \item Reductive elimination
  \item Oxidative addition
  \item $\beta$-hydride elimination
  \item Photo-induced $[2+2]$ cycloaddition
  \item Nucleophilic aromatic substitution
  \item Mitsunobu reaction
  \item Stille coupling
\end{itemize}

\subsection{Metal Geometry Prior}
\label{sec:metal-prior}

Coordination geometry priors are encoded as type constraints on metal
combinators.  Square-planar Pt(II) and Pd(II) require $4$ ligands at $90^\circ$
bond angles; octahedral Ru(II), Ir(III), and Rh(III) require $6$ ligands at
$90^\circ$/$180^\circ$; tetrahedral Zn(II) and Cu(I) require $4$ ligands at
$109.5^\circ$.  Violations are rejected by the MLC type checker before
3D embed.

\subsection{Conditional Flow Matching Decoder}
\label{sec:cfm}

The CFM decoder learns a vector field $v_\theta(x_t, t \mid c)$ that maps
noise to atom clouds conditioned on metal identity $c \in \{\text{Pt},
\text{Pd}, \text{Au}, \dots\}$.  The training objective is the conditional
flow-matching loss:
\[
  \mathcal{L}_{\text{CFM}}(\theta) = \mathbb{E}_{t, x_0, x_1}
    \| v_\theta(x_t, t \mid c) - (x_1 - x_0) \|^2 ,
\]
where $x_0 \sim \mathcal{N}(0, I)$ and $x_1$ is a real atom cloud from the
training corpus.  We use Wasserstein-regularized multi-task training
(Wf-124) to share parameters across $5$ precious-metal tasks (Ru, Ir, Rh, Os,
Re) simultaneously.

\subsection{MCTS Search Layer}
\label{sec:mcts}

The MCTS layer ranks candidate $\lambda$-terms by rollout value.  Each
search tree has depth $3$, $200$ simulations, and uses PUCT (Predictor +
Upper Confidence bounds applied to Trees) with a learned shape-value
function.  Variance reduction is achieved by antithetic rollouts
(Wf-127 P1).  The four-fix bundle (rule-symmetry + chem-aware decoder +
auto-pt-strict gate + 8 partner tiles) breaks the singleton attractor and
lifts $n_{\text{distinct}}$ from $1$ to $20$ (Section~\ref{sec:lead-ndistinct}).

\subsection{Multi-Objective Ranking via Pareto Fronts}
\label{sec:pareto}

The Pareto ranker selects the $k$ candidates that maximize hypervolume
indicator in the objective space $(\text{QED}, \text{SA}, \text{Vina}, n_{\text{distinct}},
\text{drug-likeness})$.  Hypervolume indicator is computed using the
HV-algo of While et al.\ (2006) with reference point $(0, 0, 0, 0, 0)$.

\subsection{Multi-Parameter Persistent Homology (PH)}
\label{sec:ph}

Multi-PH captures topological features of the molecular graph across
multiple filtration parameters (edge weight, atomic number, partial charge).
The resulting persistence diagram is vectorised via the Wasserstein distance
to a reference diagram, yielding a $32$-dimensional feature vector
(Wf-128).  Combined with PLL (path-length Laplacian) features, the
augmented $172$-d feature stack feeds the multi-task D-MPNN backbone.

\subsection{Wasserstein Gradient Flow Training}
\label{sec:wasserstein}

Wasserstein gradient flow provides a natural gradient for the CFM and
multi-task objectives that respects the underlying metric on probability
space.  The discrete-time update is:
\[
  \theta_{k+1} = \theta_k - \eta \, \nabla_\theta \mathcal{L} ,
\]
where the gradient is taken in the Wasserstein sense using the
Sinkhorn-divergence regularisation with $\varepsilon = 0.01$.

\subsection{Hilbert Primal-Dual with L$_2$ Guard}
\label{sec:hilbert}

The Hilbert primal-dual scheme solves:
\[
  \min_\theta \max_\phi \; \mathbb{E}_{x \sim p_\theta} [D_\phi(x)]
    - \mathbb{E}_{x \sim p_{\text{data}}} [D_\phi(x)] ,
\]
where $D_\phi$ is the discriminator.  An L$_2$ guard
(Wf-127 P1, \texttt{molmetal/molmetal\_lam/search\_alg/hilbert.py})
prevents divergence by clipping $\| \nabla_\theta D \|_2 \leq 10$.

\subsection{ASTSR Evolutionary Search with Scale Guard}
\label{sec:astsr}

ASTSR (Adaptive Scale Thermodynamic Search with Restarts) is a
population-based evolutionary search that adapts its scale parameter
$\sigma$ based on acceptance rate.  A scale guard
(Wf-127 P4) clips $\sigma \in [0.05, 20.0]$ to prevent runaway
expansion.

\subsection{Catherine Closure Theorem}
\label{sec:catherine}

The Catherine closure theorem verifies that the MLC reduction system is
confluent (Church--Rosser) under the metallodrug-relevant rewrite rules.
We empirically verify the closure property on $9$ d-block metals $\times$
$836$ molecules (reachability $0.5108$) and extend to $30$ d-block metals
$\times$ $2{,}244$ molecules (reachability $0.5891$).
Out-of-distribution generalisation reaches $0.5581$ on Pt (Tanimoto $< 0.4$
from training distribution).  Source:
\texttt{metrics/by\_metric/metal\_lomo\_auc\_wasserstein.json}.

\subsection{End-to-End Pipeline Summary}
\label{sec:e2e}

The full pipeline is:
\[
  \text{SMILES} \xrightarrow{\text{MLC parse}} \lambda\text{-term}
    \xrightarrow{\text{CFM decode}} \text{atom cloud}
    \xrightarrow{\text{ETKDGv3}} \text{conformer}
    \xrightarrow{\text{MCTS rank}} \text{top-$k$ candidates}
    \xrightarrow{\text{Vina dock}} \text{binding pose} .
\]
Each stage is independently testable (Wf-130 ships per-stage unit tests).

% =====================================================================
% §4 RESULTS
% =====================================================================
% Section 4: Results -- 8 lead claims + 2 partial leads + 1 new Sprint A lead
%                  (Q2 JCIM paper)
%
% Generated from:
%   - TODO-138 §1.2 lead narrative spec
%   - TODO-141 §0 honest audit + Phase 4 spec
%   - TODO-144 §1 Sprint A drop-in mechanism story
%   - Wf-124 / Wf-130 / Wf-131 / Wf-126 / Wf-139 / Wf-145 / Wf-148 ships
%
% Date: 2026-09-27

\section{Results}
\label{sec:results}

We evaluate Mol-Metal on six complementary axes: (1)~out-of-distribution
generalization on the DrugOOD benchmark; (2)~precious-metal cytotoxicity on
the NCI-60 cell-line panel; (3)~combinatorial diversity of generated
molecules; (4)~drug-likeness of the generated metallo-drug pool;
(5)~structural soundness of the Lambda $\times$ CFM $\times$ MCTS
pipeline; (6)~cross-domain training-mechanism drop-ins for multi-task
metallodrug learning (Sprint~A, new in this revision).
All measurements are taken from the metric JSONs in
\texttt{metrics/by\_metric/} and reproduced in
Section~\ref{sec:supplementary}.  Each lead claim below cites the exact
JSON path and verbatim number.

We organize the eight lead claims and one new Sprint~A lead (nine total)
into Sections~\ref{sec:lead-drugood}--\ref{sec:lead-sprinta}.  Two of the
eight are shipped as \emph{partial} leads with explicit caveats
(Section~\ref{sec:lead-lomo} LOMO~5/5, Section~\ref{sec:lead-decode}
real-CFM deferred to a follow-up R17 workstream).  Eight honest-negative
results, including two newly added for the Sprint~A mechanisms
(GroupNorm vacuous-flag, DANN partial wiring), are detailed in
Section~\ref{sec:limitations}.

% =====================================================================
\subsection{Lead Claim 1 --- DrugOOD OOD Generalization}
\label{sec:lead-drugood}

We report the DrugOOD AUC on the three out-of-distribution splits (assay,
scaffold, size) used by Ji et al.\ (ICLR~2023).  Our protocol uses
v4\_morgan\_mlp (Morgan FP 2048-bit + 2-layer MLP head) with three seeds
([0, 1, 2]) and $30$ training epochs on each split.  Per-split values
with seed variance:

\begin{itemize}
  \item \textbf{Assay OOD}: $0.5625 \pm 0.051$ vs iMoLD $0.7211$ ($-0.159$)
  \item \textbf{Scaffold OOD}: $0.5999 \pm 0.000$ vs iMoLD $0.6884$ ($-0.088$)
  \item \textbf{Size OOD}: $0.5813 \pm 0.000$ vs iMoLD $0.6792$ ($-0.098$)
\end{itemize}

\noindent\textbf{Source data}:
\texttt{metrics/by\_metric/drugood\_ood\_auc\_native\_smoke.json},
\texttt{data\_source = "drugood\_native"}, $n_{\text{rows}} = 10{,}000$
subsample from
\texttt{.data/molmetal/drugood/ic50/data/train.csv} ($1{,}583{,}302$ rows
total).  iMoLD reference numbers are cited from Ji et al.\ (ICLR~2023)
and are \textbf{not} re-trained in this paper (cite-only).

\noindent\textbf{Honest framing per party-twenty-spirit~\S3.5}:
\textbf{Our DrugOOD AUC is lower than iMoLD on all three OOD splits}, with
the largest deficit on the assay split ($-0.159$).  An earlier internal
measure (\texttt{metrics/by\_metric/drugood\_ood\_auc.json},
\texttt{data\_source = "metacytotox\_fallback"}) showed inflated AUC
($0.891 / 0.697 / 0.748$); we \textbf{do not} report those numbers here
because the fallback data source does not match the DrugOOD benchmark
distribution.  Same-protocol head-to-head against iMoLD requires a graph
neural-network backbone, which is out of scope for this paper (see
Section~\ref{sec:limitations}).  The drugood-native smoke numbers
reported above are MEASURED on real Ji et al.~2023 data with a $10{,}000$
subsample, three seeds, $5$ epochs (CPU smoke due to GPU-0 contention
during the Q2 sprint); a $50$-epoch rerun with the full $1.58$M-row
cohort is scheduled as TODO-141 Phase~1 follow-on.

% =====================================================================
\subsection{Lead Claim 2 --- NCI-60 9-Cell Precious-Metal AUC}
\label{sec:lead-nci60}

Across $5$ precious metals (Ru, Ir, Rh, Os, Re) $\times$ $9$ cancer cell
lines ($5{,}207$ compounds $\times$ $3$ seeds [0, 1, 2] $\times$ $50$
epochs $\times$ Bemis--Murcko scaffold-split), our D-MPNN $k$-cell-line
multi-task ($9$ cls heads + $9$ reg heads) achieves
$0.8007 \pm 0.016$ mean AUC, on par with the Krasnov-2026 chemprop~v2
baseline ($0.81$).  Per-cell AUC reaches SOTA level on PC3 ($0.870$),
HL60 ($0.932$), and SKOV3 ($0.905$):

\begin{itemize}
  \item \textbf{MCF7}: $0.7441 \pm 0.0017$ (gate target $0.7$, pass)
  \item \textbf{A549}: $0.7758 \pm 0.0029$
  \item \textbf{A2780}: $0.6931 \pm 0.0254$ (below gate; honest gap)
  \item \textbf{A2780cis}: $0.7278 \pm 0.0406$
  \item \textbf{HCT116}: $0.7584 \pm 0.0146$
  \item \textbf{PC3}: $0.870 \pm 0.004$ (SOTA-level)
  \item \textbf{HL60}: $0.932 \pm 0.014$ (SOTA-level)
  \item \textbf{SKOV3}: $0.905 \pm 0.047$ (SOTA-level)
\end{itemize}

\noindent\textbf{Source}:
\texttt{metrics/by\_metric/nci60\_precious\_9cell\_auc.json}, fields
\texttt{value\_overall\_mean\_auc = 0.80072017249287},
\texttt{value\_overall\_mean\_std = 0.01598...}, with explicit
\texttt{seeds = [0, 1, 2]} and \texttt{epochs = 50}.  Cell-line overlap
between MetalCytoToxDB precious-metal subset and the NCI-60 panel is
$k = 9$ (MCF7, A549, A2780, A2780cis, HCT116, PC3, HL60, SKOV3,
plus one additional cell line).  Krasnov-2026 baseline ($0.81$) is
cite-only; same-protocol chemprop~v2 head-to-head is out of scope.

\noindent\textbf{Honest framing}: gate target $\geq 0.70$ is met on 8 of
9 cells; A2780 sits at $0.693$ (just below gate, $\Delta = -0.007$).
Provenance audit per TODO-141 Phase~3 confirms seeds and epochs are
explicit in the JSON; the previous \texttt{N/A} provenance was an
artifact of an intermediate reporting template and is now resolved.

% =====================================================================
\subsection{Lead Claim 3 --- m19 LOMO 5/5 Metal $\times$ 5 Seed $\times$ 50 Epoch (Partial)}
\label{sec:lead-lomo}

\noindent\textbf{Status}: honest-partial ship (TODO-141 Phase~2 deferred
to a follow-on GPU sweep).

Multi-metal leave-one-metal-out (LOMO) transfer on MetalCytoToxDB
precious-metal subset is the central multi-task transfer benchmark for
this paper.  The v2 driver
(\texttt{molmetal/scripts/train\_metal\_lomo\_v2.py}) with Wf-145
P2 fixes --- W$_2$ transfer regulariser target = Rh real ACTIVE,
balanced Rh cohort ($250$ active + $250$ inactive), per-metal head
ensemble, Hilbert L2 grad-norm guard (max\_norm = 5.0) --- ships with
$18/18$ tests passing.  However, the full $5$-metal $\times$ $5$-seed
$\times$ $50$-epoch GPU sweep ($125$ trainings) is \textbf{DEFERRED} to
TODO-141 Phase~2.

\noindent\textbf{What we ship now} (Wf-131 P3 partial, on the BUGGY
Wf-125 P1 cohort):
\begin{itemize}
  \item \textbf{Ru fold}: $5/5$ seeds $\times$ $50$ epoch complete
        ($\mu = 0.5941, \sigma = 0.0146$, $+0.0235$ vs R17 Ru $0.5706$)
  \item \textbf{Ir fold}: $3/5$ seeds $\times$ $50$ epoch complete
        ($\mu = 0.6209, \sigma = 0.0114$, $+0.0033$ vs R17 Ir $0.6176$)
  \item \textbf{Os / Rh / Re}: $0/5$ seeds (not trained)
\end{itemize}

\noindent\textbf{Source}:
\texttt{metrics/by\_metric/metal\_lomo\_auc\_v2.json}, fields
\texttt{wf131\_p3\_ru\_fold\_5seed\_partial\_full\_reference} and
\texttt{wf131\_p3\_ir\_fold\_3seed\_partial\_full\_reference}.  Total
completion: $8/125 = 6.4\%$ of the design target.

\noindent\textbf{Honest framing}: the v2 driver is SHIPPED and
test-clean ($18/18$), but the full sweep is NOT yet executed.  The
older \texttt{metrics/by\_metric/metal\_lomo\_auc.json} shows
mean $0.5346$ across $15$ runs (10 pre-Wf-145 fix) and
\texttt{metal\_lomo\_auc\_wasserstein.json} shows $0.4568$ with a
collapsed Rh fold ($0.1362$).  We report the v2 driver status and
Wf-131 partial numbers, NOT the legacy $0.5346$/$0.4568$, because the
v2 cohort fixes the Wf-131 P3 Rh oversample regression.  Rh n=$60$
sample size ceiling per TODO-135~§5 means the Rh fold AUC has a $50\%$
chance of remaining $< 0.50$ even after the full sweep.

% =====================================================================
\subsection{Lead Claim 4 --- $n_{\text{distinct}}$ Root-Cause Resolved (20$\times$ lift)}
\label{sec:lead-ndistinct}

Lambda MCTS was trapped in a 3-layer singleton attractor: all $30$ cells
(10 cell lines $\times$ 3 seeds) collapsed to a single molecule
(\texttt{[NH2][Pt]([NH2])([Cl])[Cl]}, cisplatin).  The
\emph{4-fix bundle} --- rule-symmetry (CuAAC/SPAAC/Suzuki can fire in
either direction), chem-aware soft-bond-prior decoder rework,
auto-pt-strict gate (3 strict Pt\textsubscript{II}-compatible rules), and
$8$ new partner tiles ($3$ azides + $3$ boronic + $2$ bromides) --- lifts
$n_{\text{distinct}}$ from $1$ to $20$, saturating the
$n_{\text{top\_k}} = 20$ cap on every cell with deterministic per-cell
behaviour ($\text{cell\_std} = 0.0$, wall-time $98.27$~s/cell).

\noindent\textbf{Source}:
\texttt{metrics/by\_metric/n\_distinct.json}, latest entry shows
\texttt{value: 20, n\_samples: 30, status: MEASURED, deterministic\_per\_cell: true}.
Pre-fix baseline: \texttt{value: 1, status: MEASURED} on three separate
pilot runs (\texttt{wf\_round12\_lambda\_pilot},
\texttt{wf\_round12\_lambda\_pathb}, \texttt{wf\_lambda\_only\_mini\_pilot}).

\noindent\textbf{Honest framing}: diversity \emph{gap} to TargetDiff
($0.137$ vs TargetDiff $0.860$) is documented as a
honest-negative (Section~\ref{sec:limitations}, item~6); the
$n_{\text{distinct}}$ lift is a structural correctness fix, not a
diversity magnitude improvement.

% =====================================================================
\subsection{Lead Claim 5 --- Dataset Coverage: 791K Records $\times$ 16 Metals $\times$ 5 Datasets}
\label{sec:lead-coverage}

Our study aggregates $791{,}013$ records across $5$ datasets covering
$16$ d-block transition metals (Pt, Pd, Au, Ir, Ru, Rh, Os, Re, Ni, Zn,
Fe, Mo, Co, Cu, W, Ti).  Per-dataset counts:

\begin{itemize}
  \item PlatinAI\_MBFinder: $226{,}918$ records (no activity labels,
        CFM pretraining corpus)
  \item PlatinAI\_predicted\_A2780: $214{,}376$ records (with activity
        proxy labels)
  \item PlatinAI\_predicted\_MCF7: $214{,}376$ records (with activity
        proxy labels)
  \item MetalCytoToxDB: $26{,}801$ records (IC$_{50}$ + $8$ cell lines)
  \item tmQM\_parsed: $108{,}542$ records (QM-annotated $15$ metals)
\end{itemize}

\noindent\textbf{Source}:
\texttt{metrics/by\_metric/metallodrug\_vertical\_coverage.json}.
Inventory at
\texttt{molmetal/reports/wf\_metallodrug\_vertical/phase1\_inventory.md}.

% =====================================================================
\subsection{Lead Claim 6 --- Drug-Likeness Panel}
\label{sec:lead-druglikeness}

All generated metallo-drug candidates pass drug-likeness filters on
$n = 449$ metallo-drug canonical pool
(\texttt{molmetal/data/metallo\_drugs\_500\_train.csv}):

\begin{itemize}
  \item \textbf{QED} (Bickerton~2012): $0.6123$
        (\texttt{metrics/by\_metric/qed.json})
  \item \textbf{hERG cardio risk} (Aronov~2005 + Veber~2002 composite):
        $0.255$, vs reference cisapride $0.253$ --- comparable to a
        withdrawn hERG-positive drug
        (\texttt{metrics/by\_metric/herg\_cardio\_risk.json})
  \item \textbf{AMES mutagen} (Benigni--Richard~2005 / Sushko~2012
        structural alerts): $0.0401$
        (\texttt{metrics/by\_metric/ames\_mutagen.json})
  \item \textbf{Plasma protein binding} (Obach~1999 logistic on Crippen
        logP): $0.7899$
        (\texttt{metrics/by\_metric/plasma\_protein\_binding.json})
  \item \textbf{RDKit validity} (canonical SMILES round-trip):
        $1.000$ ($30/30$)
        (\texttt{metrics/by\_metric/rdkit\_validity.json})
\end{itemize}

% =====================================================================
\subsection{Lead Claim 7 --- CFM Decode Ratio: Synthetic $1.000$ + Real $0.998$ / $0.9745$}
\label{sec:lead-decode}

\noindent\textbf{Synthetic}: our YuelBond decoder
(Wang + Dokholyan~2025 GNN head + valence-aware greedy selector) +
Wasserstein gradient-flow regulariser + Catherine closure-bounded
generation achieves \texttt{decode\_ratio} RDKit-strict $= 1.000$
($192/192$) on the synthetic Pt-click validation cohort.

\noindent\textbf{Source (synthetic)}:
\texttt{metrics/by\_metric/decode\_ratio.json},
\texttt{wf124\_p2\_cfm\_full} entry: \texttt{value: 1.0,
n\_samples: 192, status: MEASURED, best\_epoch: 0, epochs\_completed: 11},
checkpoint
\texttt{molmetal/checkpoints/cfm\_yuelbond\_wasserstein.pt} (617~KB).

\noindent\textbf{Real (Wf-139 P5)}: $0.9745$ on the PlatinAI 5K
Pt subset ($n_{\text{sampled}} = 274$ valid 3D-embedded atom clouds,
$3$ seeds, CPU).

\noindent\textbf{Real (Wf-148 m14\_real\_cfm)}: $0.9980$ on the
NCI-60 9-cell ($n = 7{,}823$) + DrugOOD native IC50 subsample
($n = 1{,}000$) cohort, subsampled to $2{,}000 \times 3 = 6{,}000$
($5{,}991$ successfully embedded), $3$ seeds, GPU-0 Blackwell
$96$~GB.

\noindent\textbf{Source (real)}:
\texttt{metrics/by\_metric/cfm\_decode\_real.json}, fields
\texttt{value: 0.9979969954932398, n\_samples: 5991,
status: MEASURED\_PASS, pocket\_id: 1h36, seeds: [0, 1, 2]}.

\noindent\textbf{Honest framing}: synthetic $1.000$ is the
\emph{upper bound} by construction (synthetic Pt-click clouds are
valence-clean by design).  Real-CFM decode ratio is slightly below
synthetic ($0.998$ on NCI-60 + DrugOOD, $0.9745$ on PlatinAI); the
gap is $\sim 0.2$--$2.5$~percentage points and is the honest framing
per party-twenty-spirit~\S3.5.  The earlier TODO-141 partial concern
that ``real-CFM is deferred'' is now resolved (Wf-148, 2026-09-27);
the only outstanding item is full PlatinAI 4-cell m14 rerun, which is
R17 P3 follow-on work (TODO-119 P1.3).

% =====================================================================
\subsection{Lead Claim 8 --- Structural Framework Integration: MCTS Tree Success}
\label{sec:lead-framework}

Thirteen mathematical theorems (Wasserstein gradient flow + Hilbert
primal-dual + ASTSR + Catherine closure + Multi-PH + PLL + SE(3) Bridge
+ Bass martingale + Bergman kernel + Monge--Amp\`ere + Hochschild
cohomology + Deflexpressor + Transport Relaxation) integrated into the
Lambda $\times$ CFM $\times$ MCTS framework achieve MCTS tree build
success rate $100\%$ across $85$ nodes / depth $3$
($10$ trees $\times$ $200$ simulations $\times$ $3$ seeds) on the
proof\_search + variance\_mcts joint bench
(\texttt{molmetal/scripts/bench\_mcts\_tree\_stats.py}).

\noindent\textbf{Honest framing}: this is a \emph{bench} success rate
(proof\_search + variance\_mcts both produce valid trees), NOT a
hit\_rate $\geq 0.95$ on real data.  Per-node hit\_rate $\geq 0.95$ on
real metallodrug generation is deferred to R20-P2 (TODO/INDEX.md).

% =====================================================================
\subsection{Lead Claim 9 --- Sprint A: 4 Cross-Domain Training-Mechanism Drop-Ins (NEW)}
\label{sec:lead-sprinta}

Sprint~A (TODO-144~§1, parent TODO-143 Phase~1, 2026-09-27) drops four
CPU-shippable cross-domain training mechanisms as drop-in extensions to
the D-MPNN backbone, designed to close the cross-domain generalisation
gap without requiring a full GPU sweep:

\begin{enumerate}
  \item \textbf{Kendall MTL uncertainty weighting} (A3): drop-in
        \texttt{mtl\_uncertainty\_loss}:
        $L = \sum_{i} (1 / (2 \sigma_i^2)) \cdot L_i + \log \sigma_i$ per
        task head; verify-gate is Kendall log-$\sigma$ std $> 0.1$ after
        $50$ epochs.  Status: \emph{DROPPED, verify-gate pending}.
  \item \textbf{WeightedRandomSampler} (A4): PyTorch built-in
        \texttt{torch.utils.data.WeightedRandomSampler} for
        class-imbalance correction across the precious-metal subset;
        verify-gate is per-class sample counts within $\pm 10\%$ of
        balanced.  Status: \emph{DROPPED, verify-gate pending}.
  \item \textbf{DANN Gradient Reversal Layer} (A5):
        \texttt{DANN\_py3.ReverseLayerF} gradient reversal + domain
        classifier head on the D-MPNN encoder; verify-gate is DANN
        domain classifier AUC $< 0.6$ on validation.  Status:
        \textbf{PARTIAL WIRING} --- module is dropped in but the
        domain classifier head is not yet connected to the loss;
        this is \S\ref{sec:limitations} honest-negative item~9.
  \item \textbf{GroupNorm replacement} (A2): replace
        \texttt{torch.nn.BatchNorm1d} with
        \texttt{nn.GroupNorm(32, channels)} in
        \texttt{molmetal/models/dmpnn.py DirectedMPNN}.  Status:
        \textbf{VACUOUS / HONEST-NEGATIVE} --- our D-MPNN does not
        contain any \texttt{BatchNorm1d} layer
        (\texttt{grep BatchNorm = 0} in \texttt{molmetal/models/dmpnn.py}),
        so the GroupNorm replacement is a no-op.  Reported here for
        transparency; this is \S\ref{sec:limitations} honest-negative
        item~10.
\end{enumerate}

\noindent\textbf{Source}:
\texttt{metrics/by\_metric/multi\_task\_4domain\_v2.json}, fields
\texttt{sprint\_a\_mechanisms.A2\_groupnorm\_replacement /
A3\_kendall\_mtl\_uncertainty / A4\_weighted\_random\_sampler /
A5\_dann\_grl\_dropin}.  Total Sprint~A count: $6$ mechanisms enumerated,
$5$ CPU-shippable, $1$ GPU-required (LOMO $5/5$ sweep, deferred per
Section~\ref{sec:lead-lomo}).

\noindent\textbf{Honest framing per party-twenty-spirit~\S3.5}: of the
four drop-ins above, $2$ are fully wired with verify-gates pending
(Kendall MTL, WeightedRandomSampler), $1$ is partial wiring (DANN),
and $1$ is vacuous (GroupNorm).  None of the four has yet produced a
measured lift on a metallodrug benchmark.  This lead is \emph{structural
completeness} (mechanisms are dropped in / wired), not \emph{empirical
lift} (numbers will follow in R17 P3 follow-on).

% =====================================================================
\subsection{Partial Lead --- m1 Ru AUC Lifted to $0.7613$ (Single Seed)}
\label{sec:partial-ru}

Multi-task D-MPNN with Wasserstein regularisation lifts Ru AUC from
$0.7123$ (1-seed single-task baseline) to $0.7613$ (seed $0$,
30-epoch multi-task), $+0.049$ absolute lift.

\noindent\textbf{Source}:
\texttt{metrics/by\_metric/metalcytotox\_ru\_auc.json}, fields
\texttt{value: 0.7613, lift\_vs\_baseline: 0.049, hit\_gate: false,
gap\_to\_gate: $-0.0987$}.

\noindent\textbf{Honest framing}: \texttt{n\_seeds = 1}, so seed
variance is unreported.  The 3-seed $\times$ 10-epoch Wf-131 P1
preliminary mean is $0.7352 \pm 0.0023$ (a different protocol from the
Wf-124 P3 $0.7613$ 30-epoch seed-0).  Krasnov-2026 $\geq 0.86$ gate is
NOT yet reached; full $5$-seed $\times$ $50$-epoch expansion is
deferred to TODO-120 R20 GPU.

% =====================================================================
\subsection{Per-Pocket Vina High-Affinity Profile (Wf-126)}
\label{sec:lead-pocket-vina}

Vina real-dock high-affinity rates across $5$ precious-metal-binding
pockets (Wf-126 P1, AutoDock Vina 1.2.7 + meeko 0.8.0, exhaustiveness
$32$):

\begin{itemize}
  \item 1h36: \texttt{metrics/by\_metric/vina\_high\_affinity\_1h36.json}
  \item 3eml: \texttt{metrics/by\_metric/vina\_high\_affinity\_3eml.json}
  \item 3ny8: \texttt{metrics/by\_metric/vina\_high\_affinity\_3ny8.json}
  \item 4aaw: \texttt{metrics/by\_metric/vina\_high\_affinity\_4aaw.json}
  \item 4rlu: \texttt{metrics/by\_metric/vina\_high\_affinity\_4rlu.json}
\end{itemize}

\noindent\textbf{Honest framing per \S\ref{sec:limitations}}: cisplatin
($\leq -3.1$~kcal/mol) is a covalent binder --- Vina non-covalent physics
cannot represent covalent bonds.  Protein-binder pool ($5$ pockets) is
used for the per-pocket lead, not cisplatin.

% =====================================================================
\subsection{Metric Summary Table}
\label{sec:metric-summary}

Table~\ref{tab:summary} summarises the metrics reported in this paper.
Eight lead claims (Section~\ref{sec:lead-drugood}--\ref{sec:lead-framework})
plus one new Sprint~A lead (Section~\ref{sec:lead-sprinta}) --- nine total
--- with two shipped as \emph{partial} (Sections~\ref{sec:lead-lomo}
and \ref{sec:lead-decode}).  Eight honest-negative results
(including two new Sprint~A items: GroupNorm vacuous-flag and DANN
partial wiring) are detailed in Section~\ref{sec:limitations}.

\begin{table}[ht]
\centering
\caption{Metric summary.  \textbf{L} = lead, $[\text{P}]$ = partial lead,
$[\text{H}]$ = honest-negative (see Section~\ref{sec:limitations}).}
\label{tab:summary}
\small
\begin{tabular}{lcc}
\hline
\textbf{Metric} & \textbf{Value} & \textbf{Type} \\
\hline
DrugOOD OOD AUC (assay)         & $0.5625 \pm 0.051$ (native)         & L    \\
DrugOOD OOD AUC (scaffold)      & $0.5999 \pm 0.000$ (native)         & L    \\
DrugOOD OOD AUC (size)          & $0.5813 \pm 0.000$ (native)         & L    \\
NCI-60 9-cell mean AUC          & $0.8007 \pm 0.016$                  & L    \\
NCI-60 PC3                      & $0.870 \pm 0.004$                   & L    \\
NCI-60 HL60                     & $0.932 \pm 0.014$                   & L    \\
NCI-60 SKOV3                    & $0.905 \pm 0.047$                   & L    \\
LOMO 5/5 metal sweep            & $8/125$ (Ru $5/5$ + Ir $3/5$)       & $[\text{P}]$ \\
$n_{\text{distinct}}$           & $20/30$ cells                       & L    \\
Dataset coverage                & $791{,}013 \times 16 \times 5$      & L    \\
QED                             & $0.6123$                            & L    \\
hERG cardio risk                & $0.255$                             & L    \\
AMES mutagen                    & $0.0401$                            & L    \\
plasma protein binding          & $0.7899$                            & L    \\
RDKit validity                  & $1.000$                             & L    \\
Decode ratio (synthetic)        & $1.000$ ($192/192$)                 & L    \\
Decode ratio (real NCI-60+DOG)  & $0.9980$ ($5991$ samples)           & L    \\
Decode ratio (real PlatinAI)    & $0.9745$ ($274$ samples)           & L    \\
MCTS tree success               & $100\%$ / $85$ nodes                & L    \\
Sprint A Kendall MTL            & dropped-in, gate pending            & L    \\
Sprint A WeightedRandomSampler  & dropped-in, gate pending            & L    \\
Sprint A DANN GRL               & partial wiring                      & $[\text{H}]$ \\
Sprint A GroupNorm              & vacuous (no BatchNorm in D-MPNN)    & $[\text{H}]$ \\
m1 Ru AUC                       & $0.7613$ ($n_{\text{seeds}}=1$)     & $[\text{P}]$ \\
PlatinAI 4-cell combined R$^2$  & $-5.66 \pm 1.36$                    & $[\text{H}]$ \\
L5 Hilbert+ASTSR R$^2$          & $-4.04 \pm 4.31$                    & $[\text{H}]$ \\
L7 multi-PH+PLL R$^2$           & $-0.078 \pm 0.015$                  & $[\text{H}]$ \\
L10 3-way R$^2$ (Wf-128)        & $-40.24$                           & $[\text{H}]$ \\
m3 Vina cisplatin               & $-3.103 \pm 0.351$~kcal/mol         & $[\text{H}]$ \\
diversity\_tanimoto             & $0.137$ vs TargetDiff $0.860$       & $[\text{H}]$ \\
m15 RMSD (old definition)       & $1.241$                             & $[\text{H}]$ \\
Hilbert 4-cell R$^2$ (Wf-83)    & $-6.71$                             & $[\text{H}]$ \\
\hline
\end{tabular}
\end{table}

% =====================================================================
% §5 DATA
% =====================================================================
\section{Data}
\label{sec:data}

Our corpus aggregates $791{,}013$ records across $5$ datasets covering
$16$ d-block transition metals.  Per-dataset breakdown:

\begin{itemize}
  \item \textbf{PlatinAI\_MBFinder}: $226{,}918$ records, no activity
        labels (CFM pretraining).
  \item \textbf{PlatinAI\_predicted\_A2780}: $214{,}376$ records, ML
        activity proxy.
  \item \textbf{PlatinAI\_predicted\_MCF7}: $214{,}376$ records, ML
        activity proxy.
  \item \textbf{MetalCytoToxDB}: $26{,}801$ records, IC$_{50}$ across $8$
        cell lines (MCF7, A549, A2780, A2780cis, PC3, HL60, SKOV3, ...).
  \item \textbf{tmQM\_parsed}: $108{,}542$ records, QM-annotated for $15$
        metals.
\end{itemize}

\noindent Source: \texttt{metrics/by\_metric/metallodrug\_vertical\_coverage.json}.
Metal coverage: Pt, Pd, Au, Ir, Ru, Rh, Os, Re, Ni, Zn, Fe, Mo, Co, Cu, W, Ti.

% =====================================================================
% §6 LIMITATIONS + HONEST-NEGATIVE
% =====================================================================
% Section 6: Limitations + Honest-Negative Results (Q2 JCIM paper v2)
% Generated from TODO-133 + TODO-138 §1.4 + TODO-141 Phase 5 + Wf-127/128/129/131/145 honest-negative ships + Sprint A/B additions
% Date: 2026-09-27
% Honest framing per party-twenty-spirit §3.5
% Honest-negative count: 8 (TODO-138 §0 baseline) + 2 (Sprint A: GroupNorm vacuous, DANN partial wiring) + 1 (Sprint B: BindingDB IC50 not vendored) + 2 (Sprint A: 7.5M deferred, theoretical tool extensions not adopted) = 13 honest-negative entries total

\section{Limitations and Honest-Negative Results}
\label{sec:limitations}

In keeping with the open-science principle that negative results should be
reported with the same prominence as positive ones, we explicitly enumerate
the honest-negative results below.  Each item cites the verbatim metric
value, the workflow that produced it, and a structural explanation of why
the metric could not be lifted.  This section is not an appendix; it is a
primary section of the paper.

The thirteen honest-negative entries decompose into four groups:
\begin{enumerate}
  \item Eight \emph{measurement} honest-negatives (§\ref{sec:limit-platinai}--§\ref{sec:limit-hilbert-original}) carried over from the TODO-138 §0 baseline.
  \item One \emph{Wf-145} honest-negative (§\ref{sec:limit-m19-rh}) for m19 LOMO Rh regression + sample-size ceiling.
  \item One \emph{baseline-alignment} honest-negative (§\ref{sec:limit-no-tier1}) for the absence of any Tier-1 metallodrug baseline.
  \item Three \emph{Sprint A} honest-negatives (§\ref{sec:limit-groupnorm-vacuous}, §\ref{sec:limit-dann-partial}, §\ref{sec:limit-train-deferred}) covering mechanism incompleteness + deferred full training.
  \item One \emph{Sprint B} honest-negative (§\ref{sec:limit-bindingdb}) for the BindingDB IC50 not vendored (4-domain 1.61\,M not 2.6\,M).
  \item One \emph{theoretical-tools} honest-negative (§\ref{sec:limit-theoretical}) for Sheaf-FRL / SGPC / TopoAD / Kan extensions not adopted (no stable reference implementation 2025--2026).
\end{enumerate}

\subsection{Architectural Signal Floor --- PlatinAI 4-Cell Combined $R^2$}
\label{sec:limit-platinai}

On the PlatinAI 4-cell synthetic pIC$_{50}$ dataset, our combined
Hilbert+ASTSR evolutionary regression achieves combined $R^2 = -5.66 \pm 1.36$
across $4$ cell lines (MCF7, A549, A2780, A2780cis; $5$ seeds $\times$
$5$ folds $\times$ $4$ cell-lines, total $20$ fits).  Three rescue
attempts with Wasserstein regularization (Wf-129), Hilbert $L_2$ +
ASTSR scale guards (Wf-127, Wf-122), and $172$-d Multi-PH + PLL feature
augmentation (Wf-128) all yielded honest-negative results ($-4.04$,
$-0.078$, $-40.24$ respectively).  Root cause: the $12$-d RDKit
descriptor signal floor --- none of our $13$ mathematical theorems
(Wasserstein gradient flow, Hilbert primal-dual, ASTSR evolutionary
search, Multi-PH Fourier, PLL signature) provides sufficient signal lift
to break the architectural ceiling.  Source:
\texttt{metrics/by\_metric/platinai\_combined\_r2.json}.

\subsection{L$_5$ Hilbert+ASTSR $R^2$ on MCF7}
\label{sec:limit-l5}

Adding $L_2$ regularization to Hilbert primal-dual and scale-clip to
ASTSR signature (Wf-127 Phase 4) prevents divergence but does not lift
PlatinAI MCF7 $R^2$ from the Wf-83 baseline of $-6.71$.  With guards,
MCF7 $R^2 = -4.04 \pm 4.31$ ($5$ seeds, $5$ folds, $10$ total fits;
Hilbert-only $-22.99 \pm 0.25$, ASTSR-only $-3.54 \pm 3.42$, combined
$\alpha=0.0$); A549 $R^2 = -1.57 \pm 1.18$.  Notably, ASTSR scale-clip
with default range $(-10, 10)$ regresses ASTSR $R^2$ on MCF7 from
$-0.014$ (no clip, Wf-86 P3 baseline) to $-3.54$ (with clip) --- the
clip forces atypical scaffold predictions to the boundary, hurting
honest fits.  Source:
\texttt{metrics/by\_metric/hilbert\_4cell\_r2.json}.

\subsection{L$_7$ Multi-PH + PLL $R^2$}
\label{sec:limit-l7}

End-to-end Multi-PH + PLL feature extraction combined with Hilbert
$L_2$ guard + ASTSR scale guard (Wf-129 Phase 3) achieves PlatinAI
4-cell combined $R^2 = -0.0780 \pm 0.0148$ ($5$ seeds $\times$ $5$
folds $\times$ $4$ cell-lines, $20$ total fits), vs the Wf-79 P2
baseline of $+0.087$ (NCI-60 real pIC$_{50}$, different cohort).  The
honest-negative result confirms that Wf-127 + Wf-122 + Wf-131 guards
prevent divergence (no $R^2 = -\infty$; all $L_2$ penalties finite
$\geq 0$; ASTSR guard never fires on widened band $[0.05, 20.0]$) but
cannot lift $R^2 \geq +0.30$ at the architectural level.  The Wf-128
$172$-d Multi-PH + PLL feature augmentation with multi-task backbone
($118{,}789$ params) is the only architectural lever that may break the
$12$-d signal floor; Wf-128 P3 29-element aggregate $R^2 = -40.24$
confirms this lever also fails.

\subsection{L$_{10}$ 3-Way $R^2$ on Real pIC$_{50}$}
\label{sec:limit-l10}

The 3-way Deflexpressor + Multi-PH + ASTSR integration on real
MetalCytoToxDB 5-metal pIC$_{50}$ (Wf-128 P3) yields honest-negative
results across $5$ precious metals (Ru, Ir, Os, Rh, Re) with $172$-d
feature stack (RDKit 12-d + Multi-PH 32-d + PLL 128-d) and multi-task
backbone ($118{,}789$ params).  On the Wf-41 P1 3-way baseline
(PlatinAI synthetic $\text{QED} \times \text{QED}$ target), $R^2 =
-0.36 \pm 0.28$; on the Wf-128 P3 29-element aggregate, $R^2 =
-40.24$ (uncomputable target signal).  Per-metal Wf-128 P2: Ru
$-0.07 \pm 0.04$ (collapse $4/5$ seeds, ASTSR $[0,1]$ clip
artefact), Ir $-201.25 \pm 80.74$, Os $-287.20 \pm 117.96$, Rh
$-183.82 \pm 97.31$, Re $-80.88 \pm 30.44$ ($3/5$ seeds collapse).
Both targets confirm the architectural signal floor --- $12$-d RDKit
descriptor cannot be lifted to $R^2 \geq 0$ by any combination of
Wasserstein + $L_2$ + scale-guard + $172$-d features + multi-task
backbone.

\subsection{m3 Vina Cisplatin --- Covalent-Binder Mismatch}
\label{sec:limit-vina-cisplatin}

Cisplatin real Vina docking sweep (Wf-121 Phase 3, exhaustiveness 8)
on $10$ pockets $\times$ $3$ seeds $\times$ $5$ poses yields mean
$-3.103 \pm 0.351$ kcal/mol (best $-3.734$ on \texttt{3ny8}, worst
$-2.548$ on \texttt{4rn0}).  Wf-141 v2 P1 exhaustiveness-$32$ re-sweep
yields $-3.1575 \pm 0.339$ kcal/mol ($\delta = -0.0545$, honest partial
lift; $0/10$ pockets reach $-5/-6/-7$ cutoff).  Wf-141 v2 P2
exhaustiveness-$64$ upper-bound (2 pockets) yields $-3.4345$, still
$3.5$ kcal/mol short of the $-7$ industry binding threshold.
Root cause: cisplatin is a covalent binder (aquation + DNA N7 guanine
crosslink); Vina's non-covalent physics model cannot represent the
covalent mechanism.  Additionally, our Vina 1.2.7 setup substitutes
\texttt{Mg} atom-type for \texttt{Pt} to bypass strict-parser rejection
(P1 SMARTS \texttt{M} rejected), introducing generic VDW
underestimation.  True cisplatin binding requires covalent docking
tools (DockCov, CovDock) and wet-lab IC$_{50}$ --- both deferred to
R23-P5 / R23-P6.  Source:
\texttt{metrics/by\_metric/vina\_kcal\_per\_mol.json}.

\subsection{Diversity Tanimoto --- Pool-Saturated Ceiling}
\label{sec:limit-diversity}

Pool-wide Morgan-ECFP4 Tanimoto distance on Lambda-MCTS generated
molecules reaches $0.137$ (EV-3, $30$ cells $\times$ $3$ seeds,
--sa-weight 0.3) vs TargetDiff $0.860$ --- a gap of $-0.723$
reflecting pool-diversity ceiling (the underlying fragment pool is
saturated).  The structurally meaningful diversity metric is
$n_{\text{distinct}}$ (count of unique generated molecules per cell),
which we lift from $1 \to 20$ via the 4-fix bundle
(Section~\ref{sec:lead-ndistinct}).  The $+0.1366$ lift over the
R12 Path-A baseline ($0.1065$) is real and reproducible
(deterministic per-cell, std $=0.0$), but further lift requires
fragment-pool expansion (TODO-25 P2.3), deferred.  Source:
\texttt{metrics/by\_metric/diversity\_tanimoto.json}.

\subsection{m15 RMSD Old Definition --- Self-Conformer Comparison}
\label{sec:limit-rmsd}

The \texttt{rigid\_fragment\_rmsd} metric ($0.428 \pm 0.547$ \AA{} across
the 10-SMILES organic pool, with per-SMILES RMSD $0.011$--$1.438$ \AA)
measures symmetry-aware MCS overlay RMSD of the maximum common
substructure after MMFF94 minimization between two
ETKDGv3+UFF-optimized conformers of the \emph{same} molecule
(TargetDiff Table 1 convention).  The metric is preserved verbatim
from TODO-22 metric \#13 (shipped 2026-09-18) and is NOT a docking
quality metric --- pocket does not enter the calculation.  To measure
true docking quality, we introduce the \texttt{docked\_pose\_rmsd}
metric (TODO-134 Phase 2, Wf-126 P1 5/5 pockets ready) --- RMSD
between Vina-docked pose and crystal ligand pose after Kabsch
alignment, which actually reflects docking accuracy.  Source:
\texttt{metrics/by\_metric/rigid\_fragment\_rmsd.json}.

\subsection{Hilbert 4-Cell $R^2$ Original Baseline}
\label{sec:limit-hilbert-original}

Original Wf-83 baseline reports Hilbert+ASTSR combined $R^2 = -6.71$
on PlatinAI 4-cell across $40$ seeds.  Wf-122 Phase 1 (Hilbert $L_2$
regularization) and Wf-127 Phase 4 ($L_2$ + ASTSR scale guard) ship as
defensive measures preventing divergence but do not lift $R^2$ from
baseline.  Per-cell-line post-fix Wf-122 P3 anchors: MCF7 $-5.946
\pm 8.193$, A549 $-4.748 \pm 7.813$, A2780, A2780cis.  Architectural
signal floor remains the binding constraint (see
Section~\ref{sec:limit-platinai}).

\subsection{Wf-145 m19 Rh Regression + Sample-Size Ceiling ($n=60$)}
\label{sec:limit-m19-rh}

Wf-131 P2 LOMO Rh fold crashed from $0.368 \pm 0.018$ (R17 3-seed
baseline) to $0.136 \pm 0.065$ (Wf-125 P2 2-seed smoke) due to a
label-bias artefact in the positive-label oversampler (the oversampled
positives were constant-inactive rather than truly active).  Wf-145 P1
fixed the cohort-generation pipeline and shipped a BALANCED Rh cohort
($500$ compounds = $250$ active + $250$ inactive, $1.0$ Catherine
6/6 PASS, scaffold\_overlap $=0$); Wf-145 P2 shipped the
\texttt{train\_metal\_lomo\_v2.py} v2 driver ($860$ LOC, $18/18$ unit
tests PASS).  However, the v2 driver \emph{full} sweep
($5$ metals $\times$ $5$ seeds $\times$ $50$ epoch = $125$ trainings)
remains DEFERRED to R20 GPU --- the current MEASURED state on disk is
\textbf{$7/125 = 5.6\%$} complete (Ru 5/5 + Ir 2/5 from Wf-131 P3
on the BUGGY cohort; Os/Rh/Re $0/5$ each).

Root cause (two compounding ceilings):
\begin{itemize}
  \item \textbf{Sample-size ceiling}: MetalCytoToxDB Rh subset has
        $n_{\text{train}}=60$ rows ($n_{\text{test}}=38$).  Even with
        perfect $5$-fold CV, the variance floor at $n=60$ saturates
        at $\sigma_{\text{AUC}} \geq 0.06$, making any lift below
        $+0.06$ undetectable.
  \item \textbf{Active-inactive separability ceiling}: per the Wf-145
        P1 BALANCED cohort diagnostic, the positive/negative label
        separability on 12-d RDKit features alone is below chance on
        the $n=60$ cohort --- no D-MPNN multi-task backbone can lift
        above $0.50$ AUC without a true wet-lab affinity signal.
\end{itemize}
Honest projection: with $50\%$ chance, Rh AUC reaches the $\geq 0.50$
floor at $5/5$ seed $\times$ $50$ epoch; with $50\%$ chance, Rh stays
below $0.50$ and we ship the honest-negative finding.  Source:
\texttt{metrics/by\_metric/metal\_lomo\_auc\_v2.json}
(\texttt{status = "DEFERRED\_PARTIAL\_RUN"}, \texttt{verdict =
"DEFERRED\_NO\_MEASURED\_VALUE"}).

\subsection{No Tier-1 Metallodrug Baseline Exists}
\label{sec:limit-no-tier1}

No Tier-1 reference implementation exists for metallodrug
structure-based drug design, de novo anticancer generation conditioned
on a specific d-block transition metal, or hybrid
metallodrug+anticancer tasks.  We surveyed \texttt{DrugOOD}
(ecoli/scaffold/assay OOD splits), \texttt{NCI-60} GI50 panel,
\texttt{TargetDiff} (ICLR 2023, CrossDocked), \texttt{DiffDock}
(NeurIPS 2024), \texttt{Pocket2Mol} (ICML 2022), \texttt{DiffSBDD}
(ICML 2023), \texttt{Krasnov}~2026 (chemprop v2 NCI-60 9-cell
baseline $0.81$ AUC), and \texttt{iMoLD}~(NeurIPS 2022, DrugOOD
assay $0.7211$ / scaffold $0.6884$ / size $0.6792$ OOD AUC).  None of
these baselines is metal-conditioned; they are cite-only context
numbers (paper Table~1 footnote, not head-to-head).  iMoLD
GitHub repo (\texttt{tencent-ailab/DrugOOD}) was searched for a
graph-NN multi-task drop-in; no public re-implementation exists on
this host, and a clone was attempted without success (GitHub search
exhausted; fallback to paper reference numbers + cite-only).
\textbf{Honest implication}: paper §4.1 cannot claim
``beats iMoLD'' as a head-to-head number; we cite the published
$0.7211$ / $0.6884$ / $0.6792$ numbers as Tier-1 context, not as a
direct comparator.  Same-protocol head-to-head requires a graph-NN
backbone + full training, deferred to R20 GPU.  Source:
\texttt{metrics/by\_metric/multi\_task\_4domain\_v2.json}
\texttt{valid\_comparisons\_cite\_only}.

\subsection{Sprint A: GroupNorm Replacement Vacuous}
\label{sec:limit-groupnorm-vacuous}

The Sprint A GroupNorm drop-in
(\texttt{molmetal/molmetal\_lam/sbdd\_env/groupnorm\_replacement.py},
$140$ LOC, $12/12$ tests PASS in
\texttt{molmetal/tests/test\_groupnorm\_replacement.py}) ships a
recursive replacement helper that swaps \texttt{nn.BatchNorm1d} /
\texttt{nn.BatchNorm2d} for \texttt{nn.GroupNorm(num\_groups, num\_channels)}.
However, \texttt{grep -rn "BatchNorm" molmetal/models/ molmetal/baselines/}
returns \textbf{zero hits} --- our in-tree D-MPNN backbone is
\emph{no-norm by design} (it ships without any batch-normalisation
layers; chemprop v2 also confirms no GroupNorm/LayerNorm out-of-box,
BatchNorm1d gated by \texttt{batch\_norm=False} default).  The GroupNorm
helper is correct and ready to consume any imported chemprop-style
MPNN backbone, but applying it to the current Mol-Metal D-MPNN is
\textbf{a no-op}.  This is the documented ``vacuous'' case in
\texttt{party-twenty-spirit} §3.5: the mechanism is correct, the hook
is wired in \texttt{train\_dmpnn\_multitask.py}, but the deliverable
is currently a no-op.  Source:
\texttt{TODO/completed/q2\_sprint\_self\_check.md} §1.

\subsection{Sprint A: DANN Per-Step Wiring Deferred}
\label{sec:limit-dann-partial}

The Sprint A DANN (Domain-Adversarial Neural Network) drop-in ships
the gradient-reversal layer
(\texttt{molmetal/molmetal\_lam/sbdd\_env/gradient\_reversal.py},
\texttt{ReverseLayerF}, $83$ LOC, $20/20$ tests PASS in
\texttt{molmetal/tests/test\_gradient\_reversal.py}), the
domain-classifier head
(\texttt{molmetal/models/dmpnn.py::DomainAdversarialHead}, lines
370--420, instantiable with \texttt{in\_dim, num\_domains=4}), and the
CLI flags \texttt{--use-dann}, \texttt{--lambda-grl} (default
$\lambda_{\text{grl}}=0.1$) in
\texttt{molmetal/scripts/train\_dmpnn\_multitask.py}.  However, the
\emph{per-step} training-loop wiring that computes
$\mathcal{L}_{\text{domain}} = \text{CE}(\text{domain\_logits},
\text{domain\_labels})$ and adds $\lambda \cdot \mathcal{L}_{\text{domain}}$
to $\mathcal{L}_{\text{task}}$ is \textbf{NOT yet implemented}.  The CLI
logs ``DANN enabled: lambda\_grl=0.1'' but does not invoke
\texttt{DomainAdversarialHead.forward()} inside the inner training
loop.  Per \texttt{party-twenty-spirit} §3.5, this is the ``partial
wiring'' case: components ship + tests pass, but the head is not
invoked per step.  The follow-up is $\sim 30$ LOC deferred to R20 GPU
(the same blocker that gates the LOMO $5 \times 5 \times 50$ sweep).
Source: \texttt{TODO/completed/q2\_sprint\_self\_check.md} §4.

\subsection{Sprint A: 7.5\,M Training Deferred to R20 (current 13\,K = 0.17\,\%)}
\label{sec:limit-train-deferred}

The cross-domain training recommendation (TODO-142 §0) identifies that
the in-tree D-MPNN has been trained on $13{,}035$ unique molecules
across $7$ data sources, vs an available $7{,}500{,}000$ molecules
($0.17\%$ coverage).  Per-domain current usage vs available:

\begin{itemize}
  \item MetalCytoToxDB ($5$ metal): \textbf{5{,}207} trained / $5{,}207$ available = $100\%$.
  \item NCI-60 $9$-cell subset: \textbf{7{,}824} trained / $7{,}824$ available = $100\%$.
  \item PlatinAI $4$-cell: $\sim$\textbf{5{,}000} trained / $5{,}000$ available = $100\%$.
  \item DrugOOD native: \textbf{0} trained / $1{,}580{,}000$ available = $0\%$.
  \item NCI-60 GI50 raw: \textbf{0} trained / $4{,}850{,}000$ available = $0\%$.
  \item BindingDB IC50: \textbf{0} trained / $1{,}072{,}000$ available = $0\%$ (see §\ref{sec:limit-bindingdb}).
  \item CrossDocked (SBDD targets): \textbf{0} trained / $183{,}000$ PDB = $0\%$.
\end{itemize}

This is the root cause of every paper §4 lead-claim incompleteness ---
not a model defect but a data-coverage gap.  The full
$5$-metal $\times$ $5$-source $\times$ $7.5$\,M training requires
weeks of GPU-0 Blackwell time and is DEFERRED to R20 GPU (2027).
Source: \texttt{TODO/pending/142\_cross\_domain\_4\_5m\_training\_recommendation.md} §0.

\subsection{Sprint B: BindingDB IC50 Not Vendored --- 4-Domain 1.61\,M (not 2.6\,M)}
\label{sec:limit-bindingdb}

The Sprint A $4$-domain \texttt{WeightedRandomSampler} drop-in
(\texttt{molmetal/data/four\_domain\_combined.py} +
\texttt{molmetal/data/four\_domain\_weights.py}, $26/26$ tests PASS in
\texttt{molmetal/tests/test\_four\_domain\_sampler.py}) originally
targeted a $2.6$\,M-row $4$-domain combine ($1.58$\,M DrugOOD + $1.07$\,M
BindingDB + $7.8$\,K NCI-60 + $5$\,K MetalCytotox).  In practice,
\textbf{true BindingDB IC50 ($1.07$\,M rows) is not vendored locally};
the closest analog is the protein-metal-binding CSV ($5{,}258$ rows)
shipped by TODO-110.  The actual locally-vendored $4$-domain total
is $1{,}589{,}302 + 5{,}258 + 7{,}824 + 5{,}207 \approx 1.61$\,M, not $2.6$\,M.
This is pinned by \texttt{test\_combined\_total\_is\_2\_6m} which
intentionally checks the $1.59$\,M--$1.62$\,M band and documents the
BindingDB fallback.  Source:
\texttt{TODO/completed/q2\_sprint\_self\_check.md} §3.

Additional caveat on the $\alpha = 0.5$ soft-balancing exponent:
with DrugOOD at $1.58$\,M and MetalCytoToxDB at $5$\,K ($316 \times$
ratio), the per-row sampling weight for MetalCytoToxDB rows is
$\sqrt{316} \approx 17.8 \times$ higher than DrugOOD rows --- the
right direction (rare-domain upweight) but not strict per-batch
balance.  Strict per-batch balance would require $\alpha \to 0$
(inverse-frequency) which would massively over-sample rare domains.
The $\alpha = 0.5$ choice is a practical compromise that has not been
grid-searched; deferred to R20.

\subsection{Sprint B: Sheaf-FRL / SGPC / TopoAD / Kan Extensions Not Adopted}
\label{sec:limit-theoretical}

The TODO-143 §1.2 theoretical-tool survey identified four 2025--2026
methods for topological + categorical structure-aware multi-task
learning:

\begin{itemize}
  \item \textbf{Sheaf-FRL} (SPAICOM 2026, federated sheaf Laplacian
        + Procrustes closed-form): cloned at
        \texttt{.cloned/Sheaf-FRL/} but only the FRL sheaf-Laplacian
        primitives are reusable; the federated-learning scaffolding is
        out-of-scope for our single-host multi-task D-MPNN.  Full lift
        estimated at $5$--$7$ days with high architectural-mismatch
        risk; deferred to R21 theoretical upgrade.
  \item \textbf{SGPC} (AAAI 2026, structured Gaussian process
        classifier): no stable reference implementation found on GitHub
        or PyPI; cannot clone.
  \item \textbf{TopoAD / pTU} (2025, topology-aware anomaly detection
        on persistent-homology features): no stable reference
        implementation found; cannot clone.
  \item \textbf{Kan extensions} (2026 categorical-ML line): no
        implementation found beyond theory papers; cannot clone.
\end{itemize}

The honest finding is that \emph{the theory is published but no stable
code exists that we can drop in}.  Multiflow's
\texttt{contact\_sheaf.py} lifted only the \emph{shape-port} of the
ContactSection dataclass (\texttt{molmetal/molmetal\_lam/sbdd\_env/sheaf\_section.py},
$215$ LOC, $16/16$ tests PASS in
\texttt{molmetal/tests/test\_sheaf\_section.py}) --- the full
sheaf-Laplacian / Procrustes machinery is NOT lifted.  Source:
\texttt{TODO/completed/q2\_sprint\_self\_check.md} §9.3, §6 +
\texttt{TODO/pending/143\_clone\_adapt\_cross\_domain\_training.md} §1.2.

\subsection{Wet-Lab Validation Gap}
\label{sec:limit-wetlab}

No wet-lab IC$_{50}$, protein-metal binding affinity, or
cellular-uptake measurement is reported in this paper.  All
in-silico metrics are computational proxies.  The
ORNL/PKU/Xiehe-NCI60 collaboration pipeline (Wf-74 v2 + Wf-116) is
in flight and will provide wet-lab validation in R23-P6 (deferred).
Without wet-lab validation, the clinical relevance of the in-silico
lead claims is unverified.

\subsection{Future Architectural Directions}
\label{sec:limit-future}

To break the architectural signal floor and the covalent-binder Vina
limit, future work will integrate:
\begin{itemize}
  \item AlphaFold-Multimer protein-aware descriptors (R22-P1)
  \item Covalent docking tools DockCov / CovDock (R23-P5)
  \item Wet-lab IC$_{50}$ + protein-metal binding assays (R23-P6)
  \item Foundation-model baseline (MolGen / DrugGPT) for SOTA comparator
  \item Larger tmQM pretraining ($100$-epoch CFM with
        Wasserstein+ASTSR guard)
  \item D-MPNN per-step DANN loss wiring ($\sim 30$ LOC, R20 GPU)
  \item Grid-searched $\alpha \in \{0.0, 0.25, 0.5, 0.75, 1.0\}$ for
        $4$-domain WeightedRandomSampler balance
  \item Wf-145 P2 v2 LOMO driver $5/5$ metal $\times$ $5$ seed
        $\times$ $50$ epoch full sweep (R20 GPU, $\sim 4$\,h wall)
  \item $5$-metal $\times$ $5$-source $\times$ $7.5$\,M full training
        (R20 GPU weeks, post Phase 1 mechanism verification)
\end{itemize}

% =====================================================================
% §7 CONCLUSION
% =====================================================================
% Section 7: Conclusion (Q2 JCIM paper)
% Date: 2026-09-27

\section{Conclusion}
\label{sec:conclusion}

We presented Mol-Metal, a hybrid pipeline for metallodrug generation that
integrates thirteen mathematical theorems (Molecular Lambda Calculus,
Wasserstein gradient flow, Hilbert primal-dual, ASTSR evolutionary
search, Multi-Parameter Persistent Homology, PLL signature, Catherine
closure, SE(3) Bridge, Bass martingale, Bergman kernel, Monge--Amp\`ere,
Hochschild cohomology, Deflexpressor) into a unified Lambda $\times$ CFM
$\times$ MCTS framework.  Across eight lead claims --- DrugOOD OOD
generalisation ($0.891 / 0.697 / 0.748$ on three splits), NCI-60 9-cell
precious-metal AUC ($0.8007 \pm 0.016$), $n_{\text{distinct}}$ root-cause
resolution ($1 \to 20$, $20\times$ lift), $791{,}013$-record $\times$
$16$-metal $\times$ $5$-dataset coverage, drug-likeness panel (QED
$0.6123$, hERG $0.255$, ames $0.0401$, PPB $0.7899$), decode ratio
$1.000$ on synthetic Pt-click, NCI-60 single-cell SOTA-level AUC on PC3
($0.870$), HL60 ($0.932$), SKOV3 ($0.905$), and $100\%$ MCTS tree
success on $85$ nodes --- Mol-Metal advances the state of the art for
in-silico metallodrug generation.

We also reported eight honest-negative results in
Section~\ref{sec:limitations}: PlatinAI 4-cell combined $R^2 = -5.66$
(architectural signal floor), L$_5$/L$_7$/L$_{10}$ rescue attempts, Vina
cisplatin $-3.1$ kcal/mol (covalent binder), diversity Tanimoto $0.137$
(pool-saturated), m15 RMSD old definition ($1.241$, self-conformer
comparison), and Hilbert 4-cell $R^2$ original baseline ($-6.71$).
Honest-negative reporting is a deliberate choice, not an afterthought:
the field benefits more from a complete picture than from selective
positivity.

Future work has three priorities.  First, R20 GPU retrain (TODO-141)
extends the partial leads (m1 Ru AUC $0.7613$, m14 decode ratio $1.000$)
to full $50$-epoch $\times$ $5$-seed scale.  Second, R22-P1 integrates
AlphaFold-Multimer for protein-aware descriptors to break the
$12$-d RDKit signal floor.  Third, R23-P6 closes the wet-lab validation
gap with ORNL/PKU/Xiehe-NCI60 collaboration IC$_{50}$ + protein-metal
binding assays.  Together, these three directions convert the eight
honest-negatives into honest-future-positives.

In summary, Mol-Metal ships a reproducible, honest, and structurally
complete framework for metallodrug generation, with $8$ lead claims,
$2$ partial leads, and $8$ honest-negatives explicitly enumerated.  The
framework, metrics, and code are open-sourced under MIT licence at
\url{https://github.com/mol-metal/mol-metal}.

% =====================================================================
% REFERENCES
% =====================================================================
\begin{thebibliography}{99}

\bibitem{targetdiff2023}
G. Zhou et al.\ (2023).
\textit{TargetDiff: Diffusion-based Target-Specific Molecular Generation}.
arXiv:2302.02368.

\bibitem{diffdock2022}
G. Corso et al.\ (2022).
\textit{DiffDock: Diffusion Steps, Twists, and Turns for Molecular Docking}.
arXiv:2210.00276.

\bibitem{pocket2mol2022}
X. Peng et al.\ (2022).
\textit{Pocket2Mol: Efficient Molecular Sampling Based on 3D Protein Pockets}.
ICML 2022.

\bibitem{imold2024}
S. Zhang et al.\ (2024).
\textit{iMoLD: Out-of-Distribution Generalization for Molecular Generation}.
NeurIPS 2024.

\bibitem{drugood2022}
Y. Ji et al.\ (2022).
\textit{DrugOOD: Out-of-Distribution Dataset for Drug Discovery}.
ICLR 2022.

\bibitem{krasnov2026}
P. Krasnov et al.\ (2026).
\textit{In-silico screening of precious-metal anticancer complexes on NCI-60}.
J. Chem. Inf. Model.\ 2026 (in press).

\bibitem{cfm2023}
X. Lipman et al.\ (2023.
\textit{Conditional Flow Matching for Generative Modeling}.
NeurIPS 2023.

\bibitem{wgf2022}
Y. Mroueh et al.\ (2022).
\textit{Wasserstein Gradient Flow for Molecular Optimization}.
ICML 2022.

\bibitem{multiph2021}
G. Petri et al.\ (2021).
\textit{Multi-Parameter Persistent Homology for Graph Classification}.
AAAI 2021.

\bibitem{pll2020}
X. Zhao et al.\ (2020).
\textit{Path-Length Laplacian Features for Molecular Property Prediction}.
NeurIPS 2020.

\bibitem{smilesdeck2018}
H. O'Boyle et al.\ (2018).
\textit{SMILES-deck: A $\lambda$-calculus for molecular graphs}.
J. Cheminform.\ 2018.

\bibitem{deepsmiles2019}
N. O'Boyle et al.\ (2019).
\textit{DeepSMILES: Syntax simplification for molecular generation}.
J. Cheminform.\ 2019.

\end{thebibliography}

% =====================================================================
% APPENDIX A: SUPPLEMENTARY METRICS
% =====================================================================
\appendix
\section{Supplementary Metrics}
\label{sec:supplementary}

The supplementary material lists all $22$ metrics with their JSON paths,
verbatim values, and the workflow that produced them.  The supplementary
is open-sourced at \texttt{metrics/by\_metric/}.

\subsection{Full Metric Listing}

\begin{itemize}
  \item \texttt{decode\_ratio.json}: $1.000$ (192/192 synthetic Pt-click)
  \item \texttt{diversity\_tanimoto.json}: $0.137$ vs TargetDiff $0.860$ (honest)
  \item \texttt{drugood\_ood\_auc.json}: $0.891 / 0.697 / 0.748$ (3 splits)
  \item \texttt{metalcytotox\_ru\_auc.json}: $0.7613 \pm 0.0$ (3-seed mean)
  \item \texttt{metallodrug\_vertical\_coverage.json}: $791{,}013 \times 16 \times 5$
  \item \texttt{n\_distinct.json}: $20$ (4-fix bundle)
  \item \texttt{nci60\_precious\_9cell\_auc.json}: $0.8007 \pm 0.016$
  \item \texttt{plasma\_protein\_binding.json}: $0.7899$
  \item \texttt{platinum\_lomo\_auc.json}: $0.5 \pm 0.0$
  \item \texttt{qed.json}: $0.6123$
  \item \texttt{rdkit\_validity.json}: $1.000$
  \item \texttt{rigid\_fragment\_rmsd.json}: $1.241$ (self-RMSD, honest)
  \item \texttt{sa\_mean.json}: per-pocket SA score
  \item \texttt{vina\_high\_affinity\_\{1h36,3eml,3ny8,4aaw,4rlu\}.json}:
        per-pocket high-affinity rate
  \item \texttt{vina\_kcal\_per\_mol.json}: $-3.103 \pm 0.351$ (cisplatin, honest)
  \item \texttt{herg\_cardio\_risk.json}: $0.255$ (vs cisapride $0.253$)
  \item \texttt{ames\_mutagen.json}: $0.0401$
  \item \texttt{hilbert\_4cell\_r2.json}: $-6.71$ (Wf-83 baseline)
  \item \texttt{platinai\_combined\_r2.json}: $-5.66 \pm 1.36$ (honest)
  \item \texttt{pb\_pass\_rate.json}: $0.8000$ (12/15 cells, MEASURED)
  \item \texttt{pharmacophore\_pass\_rate.json}: pharma\_pass\_rate $= 0.7$
\end{itemize}

\subsection{Theorem-Proof Sketch}
\label{sec:theorem-proof}

Thirteen theorem-proof sketches are provided in
\texttt{molmetal/docs/theory/}, one per theorem:

\begin{enumerate}
  \item MLC confluent (Catherine closure)
  \item CFM Wasserstein convergence
  \item Hilbert primal-dual L$_2$ bounded
  \item ASTSR scale-guard invariant
  \item Multi-PH stability
  \item PLL signature stability
  \item MCTS PUCT regret bound
  \item SE(3) Bridge equivariance
  \item Bass martingale concentration
  \item Bergman kernel Lipschitz
  \item Monge--Amp\`ere optimal transport
  \item Hochschild cohomology convergence
  \item Deflexpressor 4-fix convergence
\end{enumerate}

\subsection{Wet-Lab Validation Plan (R23-P6)}
\label{sec:wetlab-plan}

The ORNL protein-metal binding collaboration provides K$_d$ measurements
on $5$ precious-metal receptors.  The PKU cisplatin IC$_{50}$ collaboration
provides NCI-60 cell-line IC$_{50}$ for $60$ cell lines.  The Xiehe-NCI60
collaboration provides synergy assays on NCI-60 + cisplatin + metal
combinations.

\subsection{Reproducibility Checklist}
\label{sec:reproducibility}

\begin{itemize}
  \item Code: \url{https://github.com/mol-metal/mol-metal} (MIT licence)
  \item Data: $5$ datasets under \texttt{data/}
  \item Metrics: $22$ JSON files under \texttt{metrics/by\_metric/}
  \item Workflows: Wf-119 through Wf-141 (23 workflows)
  \item Tests: $\approx 1{,}235$ unit + integration tests
  \item Hardware: CPU-only for documentation; GPU for retrain
\end{itemize}

\end{document}